\section{Conclusion}
\label{sec:conclusion}

We reported four flat rank-$k$ curves across training regimes, a
three-seed decoupling of accuracy from final effective rank, a
linear probe that shows a matrix-CODI bottleneck encoding
\emph{less} target-predictive information than the uncompressed
pretrained hidden state, four positive-control readouts designed to
falsify the hypothesis that readout linearity causes
rank-blindness, and a negative control on vanilla GPT-2 SFT that
demonstrates the rank-$k$ probe alone is uninformative on a model
with no matrix bottleneck. Even readouts with non-constant Jacobians
(bilinear-plus-GELU, SVD-augmented, quadratic) produce flat rank-$k$
curves: the matrix-bottleneck training objective is rank-blind through
the full chain rule, not only through the final projection. The
model-level distinguisher between rank-blindness and probe-level
position-irrelevance is the three-seed decoupling, where the same loss
lands at materially different effective ranks across seeds.

The result sits downstream of the empirical observation by
\citet{rizvi2026illusion} that fine-tuned latent reasoners do not
use their latents, and it sits parallel to the February 2026
descriptive rank measurements of \citet{nazari2026rank} and
\citet{anonymous2026staterank}. We contribute a falsifiable
mechanism claim about the training objective, the positive-control
test that adjudicates it, and a negative-control methodology for
rank-style probes that future intervention studies should adopt.
Rank-$k$ ablation should not be used as a probe for superposition in
CODI-trained matrix latents; the probe does not measure what the name
suggests it measures.
